% GENERATED FILE. Edit canonical lessons, then run npm run generate:books.
% canonical-source-hash: fnv1a64:da9de04b930fcb58
% canonical-lessons: KA-C542-malige, KA-C542-hostilu, KA-C542-hajara, KA-C542-nelamalige, KA-C542-kapatu

\chapter{Rooms of a House}
\label{ch:ka-rooms-of-a-house}
\hlchaptermodality{542}

\begin{chapteropening}
\textbf{By the end of this chapter:} \emph{I can say an upper storey, a threshold, a hall, a basement and a cabinet.}

Complete the last lesson of chapter 542: I can say an upper storey, a threshold, a hall, a basement and a cabinet.
\end{chapteropening}

\section[\texorpdfstring{\kn{ಮಾಳಿಗೆ} (māḷige)}{ಮಾಳಿಗೆ (māḷige)}]{\kn{ಮಾಳಿಗೆ} (māḷige) — an upper storey; a flat roof}

\label{lesson:KA-C542-malige}



\begin{quote}
Before the new one: say the Kannada for archery, then the Kannada for to cheer on.
\end{quote}

\subsection*{You'll want to know: \kn{ಮಾಳಿಗೆ}}
\textbf{\kn{ಮಾಳಿಗೆ}} — \emph{māḷige} — \textquotedblleft{}an upper storey; a flat roof\textquotedblright{}.

\begin{quote}
\kn{ಮಕ್ಕಳು} \kn{ಮಾಳಿಗೆಯ} \kn{ಮೇಲೆ} \kn{ಗಾಳಿಪಟ} \kn{ಹಾರಿಸಿದರು}.
\end{quote}

\begin{quote}
\emph{makkaḷu māḷigeya mēle gāḷipaṭa hārisidaru.}
\end{quote}

\begin{quote}
The children flew kites on the roof.
\end{quote}

\kn{ಮಾಳಿಗೆ} \kn{ಮನೆ}, a house with an upper floor.

\begin{cousinweb}
\textbf{\kn{ಮಾಳಿಗೆ}} (\emph{māḷige}) is native Dravidian: DEDR 4796 lists Kannada \emph{māḍa}, \emph{māḷa} and \emph{māḷige}, \textquotedblleft{}a large building, an upstair house\textquotedblright{}.
\end{cousinweb}

\begin{grammarlens}[title={What you've built}]
One more word for describing what you lived and felt.
\end{grammarlens}

\subsection*{Your turn}
\begin{itemize}
\raggedright
  \item \emph{Say it:} \emph{māḷige}
  \item \emph{Say it:} \emph{māḷige}, once more
  \item \emph{Say it:} \emph{makkaḷu māḷigeya mēle gāḷipaṭa hārisidaru.} — The children flew kites on the roof.
  \item \emph{Recall:} say the Kannada for archery, then the Kannada for to cheer on, then say \emph{māḷige} again
\end{itemize}

\begin{tcolorbox}[breakable,colback=teal!4,colframe=teal!35!black,title={Before you move on}]
What does \kn{ಮಾಳಿಗೆ} mean? (\textquotedblleft{}An upper storey; a flat roof\textquotedblright{}.) Say it once more.
\end{tcolorbox}
\section[\texorpdfstring{\kn{ಹೊಸ್ತಿಲು} (hostilu)}{ಹೊಸ್ತಿಲು (hostilu)}]{\kn{ಹೊಸ್ತಿಲು} (hostilu) — a threshold; a doorstep}

\label{lesson:KA-C542-hostilu}



\begin{quote}
Before the new one: say the Kannada for to cheer on, then the Kannada for an upper storey.
\end{quote}

\subsection*{You'll want to know: \kn{ಹೊಸ್ತಿಲು}}
\textbf{\kn{ಹೊಸ್ತಿಲು}} — \emph{hostilu} — \textquotedblleft{}a threshold; a doorstep\textquotedblright{}.

\begin{quote}
\kn{ಅತಿಥಿ} \kn{ಹೊಸ್ತಿಲು} \kn{ದಾಟಿ} \kn{ಒಳಗೆ} \kn{ಬಂದರು}.
\end{quote}

\begin{quote}
\emph{atithi hostilu dāṭi oḷage bandaru.}
\end{quote}

\begin{quote}
The guest stepped over the threshold and came in.
\end{quote}

In many homes a rangoli is drawn before it each morning.

\begin{cousinweb}
\textbf{\kn{ಹೊಸ್ತಿಲು}} (\emph{hostilu}) is a native Kannada word for the threshold of a door; its origin is not settled.
\end{cousinweb}

\begin{grammarlens}[title={What you've built}]
One more word for describing what you lived and felt.
\end{grammarlens}

\subsection*{Your turn}
\begin{itemize}
\raggedright
  \item \emph{Say it:} \emph{hostilu}
  \item \emph{Say it:} \emph{hostilu}, once more
  \item \emph{Say it:} \emph{atithi hostilu dāṭi oḷage bandaru.} — The guest stepped over the threshold and came in.
  \item \emph{Recall:} say the Kannada for to cheer on, then the Kannada for an upper storey, then say \emph{hostilu} again
\end{itemize}

\begin{tcolorbox}[breakable,colback=teal!4,colframe=teal!35!black,title={Before you move on}]
What does \kn{ಹೊಸ್ತಿಲು} mean? (\textquotedblleft{}A threshold; a doorstep\textquotedblright{}.) Say it once more.
\end{tcolorbox}
\section[\texorpdfstring{\kn{ಹಜಾರ} (hajāra)}{ಹಜಾರ (hajāra)}]{\kn{ಹಜಾರ} (hajāra) — a hall; a central room}

\label{lesson:KA-C542-hajara}



\begin{quote}
Before the new one: say the Kannada for an upper storey, then the Kannada for a threshold.
\end{quote}

\subsection*{You'll want to know: \kn{ಹಜಾರ}}
\textbf{\kn{ಹಜಾರ}} — \emph{hajāra} — \textquotedblleft{}a hall; a central room\textquotedblright{}.

\begin{quote}
\kn{ಮನೆಯ} \kn{ಹಜಾರದಲ್ಲಿ} \kn{ಎಲ್ಲರೂ} \kn{ಕುಳಿತು} \kn{ಮಾತಾಡಿದರು}.
\end{quote}

\begin{quote}
\emph{maneya hajāradalli ellarū kuḷitu mātāḍidaru.}
\end{quote}

\begin{quote}
Everyone sat and talked in the hall of the house.
\end{quote}

The open room inside the front door.

\begin{cousinweb}
\textbf{\kn{ಹಜಾರ}} (\emph{hajāra}) is a Kannada word for the hall of a house; its origin is not settled.
\end{cousinweb}

\begin{grammarlens}[title={What you've built}]
One more word for describing what you lived and felt.
\end{grammarlens}

\subsection*{Your turn}
\begin{itemize}
\raggedright
  \item \emph{Say it:} \emph{hajāra}
  \item \emph{Say it:} \emph{hajāra}, once more
  \item \emph{Say it:} \emph{maneya hajāradalli ellarū kuḷitu mātāḍidaru.} — Everyone sat and talked in the hall of the house.
  \item \emph{Recall:} say the Kannada for an upper storey, then the Kannada for a threshold, then say \emph{hajāra} again
\end{itemize}

\begin{tcolorbox}[breakable,colback=teal!4,colframe=teal!35!black,title={Before you move on}]
What does \kn{ಹಜಾರ} mean? (\textquotedblleft{}A hall; a central room\textquotedblright{}.) Say it once more.
\end{tcolorbox}
\section[\texorpdfstring{\kn{ನೆಲಮಾಳಿಗೆ} (nelamāḷige)}{ನೆಲಮಾಳಿಗೆ (nelamāḷige)}]{\kn{ನೆಲಮಾಳಿಗೆ} (nelamāḷige) — a basement; a cellar}

\label{lesson:KA-C542-nelamalige}



\begin{quote}
Before the new one: say the Kannada for a threshold, then the Kannada for a hall.
\end{quote}

\subsection*{You'll want to know: \kn{ನೆಲಮಾಳಿಗೆ}}
\textbf{\kn{ನೆಲಮಾಳಿಗೆ}} — \emph{nelamāḷige} — \textquotedblleft{}a basement; a cellar\textquotedblright{}.

\begin{quote}
\kn{ಹಳೆಯ} \kn{ಸಾಮಾನುಗಳು} \kn{ನೆಲಮಾಳಿಗೆಯಲ್ಲಿ} \kn{ಇವೆ}.
\end{quote}

\begin{quote}
\emph{haḷeya sāmānugaḷu nelamāḷigeyalli ive.}
\end{quote}

\begin{quote}
The old things are in the basement.
\end{quote}

\kn{ನೆಲ}, the ground + \kn{ಮಾಳಿಗೆ}, a storey: the storey under the ground.

\begin{cousinweb}
\textbf{\kn{ನೆಲಮಾಳಿಗೆ}} (\emph{nelamāḷige}) is a native Kannada compound: \emph{nela}, \textquotedblleft{}the ground\textquotedblright{}, and \emph{māḷige}, \textquotedblleft{}a storey\textquotedblright{} (DEDR 4796).
\end{cousinweb}

\begin{grammarlens}[title={What you've built}]
One more word for describing what you lived and felt.
\end{grammarlens}

\subsection*{Your turn}
\begin{itemize}
\raggedright
  \item \emph{Say it:} \emph{nelamāḷige}
  \item \emph{Say it:} \emph{nelamāḷige}, once more
  \item \emph{Say it:} \emph{haḷeya sāmānugaḷu nelamāḷigeyalli ive.} — The old things are in the basement.
  \item \emph{Recall:} say the Kannada for a threshold, then the Kannada for a hall, then say \emph{nelamāḷige} again
\end{itemize}

\begin{tcolorbox}[breakable,colback=teal!4,colframe=teal!35!black,title={Before you move on}]
What does \kn{ನೆಲಮಾಳಿಗೆ} mean? (\textquotedblleft{}A basement; a cellar\textquotedblright{}.) Say it once more.
\end{tcolorbox}
\section[\texorpdfstring{\kn{ಕಪಾಟು} (kapāṭu)}{ಕಪಾಟು (kapāṭu)}]{\kn{ಕಪಾಟು} (kapāṭu) — a cabinet; a shelved cupboard}

\label{lesson:KA-C542-kapatu}



\begin{quote}
Before the new one: say the Kannada for a hall, then the Kannada for a basement.
\end{quote}

\subsection*{You'll want to know: \kn{ಕಪಾಟು}}
\textbf{\kn{ಕಪಾಟು}} — \emph{kapāṭu} — \textquotedblleft{}a cabinet; a shelved cupboard\textquotedblright{}.

\begin{quote}
\kn{ಪುಸ್ತಕಗಳು} \kn{ಗಾಜಿನ} \kn{ಕಪಾಟಿನಲ್ಲಿ} \kn{ಇವೆ}.
\end{quote}

\begin{quote}
\emph{pustakagaḷu gājina kapāṭinalli ive.}
\end{quote}

\begin{quote}
The books are in the glass cabinet.
\end{quote}

\kn{ಬೀರು} is a big wardrobe; a \kn{ಕಪಾಟು} is often set into the wall.

\begin{cousinweb}
\textbf{\kn{ಕಪಾಟು}} (\emph{kapāṭu}) is usually traced to Sanskrit \emph{kapāṭa}, \textquotedblleft{}a door, the panel of a door\textquotedblright{} (Monier-Williams): a cupboard is what has the doors. The path is not settled.
\end{cousinweb}

\begin{grammarlens}[title={What you've built}]
One more word for describing what you lived and felt.
\end{grammarlens}

\subsection*{Your turn}
\begin{itemize}
\raggedright
  \item \emph{Say it:} \emph{kapāṭu}
  \item \emph{Say it:} \emph{kapāṭu}, once more
  \item \emph{Say it:} \emph{pustakagaḷu gājina kapāṭinalli ive.} — The books are in the glass cabinet.
  \item \emph{Recall:} say the Kannada for a hall, then the Kannada for a basement, then say \emph{kapāṭu} again
\end{itemize}

\begin{tcolorbox}[breakable,colback=teal!4,colframe=teal!35!black,title={Before you move on}]
What does \kn{ಕಪಾಟು} mean? (\textquotedblleft{}A cabinet; a shelved cupboard\textquotedblright{}.) Say it once more.
\end{tcolorbox}
